\documentclass[12pt, a4paper]{article} % Classe: article, report o book [8]

% --- Lingua e Codifica (Supporto Italiano) ---
\usepackage[utf8]{inputenc}  % Codifica input
\usepackage[T1]{fontenc}      % Codifica font
\usepackage[italian]{babel}   % Lingua italiana [2]

% --- Impaginazione e Margini ---
\usepackage[margin=2.5cm]{geometry} % Imposta margini (2.5cm su tutti i lati)

% --- Pacchetti Scientifici/Matematici ---
\usepackage{amsmath, amssymb, amsfonts, amsthm} % Matematica avanzata [2]
\usepackage{graphicx} % Per inserire immagini
\usepackage{hyperref} % Per collegamenti ipertestuali
\usepackage{booktabs} % Per tabelle professionali

% --- Formattazione Testo e Altro ---
\usepackage[document]{ragged2e}
\usepackage{xcolor}    % Per usare i colori 


\title{Stefania Benazzi}
\date{30-12-2025}

\begin{document}

\maketitle

\section{Ruolo}
Stefania Benazzi è una responsabile del reparto logistica con esperienza di più di 10 anni all'interno dell'azienda e si occupa di :
\begin{itemize}
    \item Coordinamento delle attività di spedizione e ricezione merci.
    \item Gestione delle relazioni con i fornitori di trasporto e spedizionieri.
    \item Supervisione della documentazione di trasporto e conformità doganale.
    \item Ottimizzazione dei processi logistici per migliorare l'efficienza operativa.
    \item Coordinamento con altri reparti (commerciale, magazzino, produzione) per garantire una gestione fluida delle commesse.
\end{itemize}

\section{Usi dell'attuale Access:}
Stefania per quanto indicata come non utilizzatrice di Access , in realtà lo utilizza marginalmente per:
\begin{itemize}
    \item Controllare i dati di trasporto e le specifiche richieste per le spedizioni.
    \item Correzione degli Incoterms e degli indirizzi di spedizione in caso di errori.
    \item Verifica delle date di consegna previste.
\end{itemize}

\section{Gestione attuale e flussi di lavoro}

\subsection{Il "File Logistico"}
L'intero reparto logistica ruota attorno a un file Excel centralizzato che
raccoglie tutti gli ordini emessi dai commerciali che vanna gestiti all'interno
del flusso di lavoro aziendale. Le funzioni principali includono:
\begin{itemize}
    \item Monitoraggio delle date previste di spedizione.
    \item Identificazione delle necessità operative: montaggio, programmazione PLC,
          dimensioni dei silos.
    \item Gestione dei permessi eccezionali.
    \item Distribuzione delle informazioni ai reparti specifici (es. montatori, reparto
          elettrico/PLC).
\end{itemize}
A seguito dell'inserimento di un ordine , il reparto logistaca verifica che tutti i documneti utili siano salvati in cartella di rete e procede con la pianificazione delle spedizioni, richiesta permessi.
\textbf{(Da fari dare il File)}.
La comunicazione che è stata create una commessa o RA avviene tramite email inviata dal commerciale al reparto logistica alla  mail comune dell'intero reparto.

\subsection{Gestione documentale e cartaceo}
Nonostante la digitalizzazione, il cartaceo rimane indispensabile nell'attuale
organizzazione ("Togliere il cartaceo in questo momento in Cepi non è la cosa
migliore").
\begin{itemize}
    \item \textbf{Cartella S: e RA:} I documenti vengono salvati in cartelle di rete specifiche. Le copie tecniche delle RA (Richieste d'Acquisto/Commesse) vengono stampate.
    \item \textbf{RA Miste:} Per commesse sia meccaniche che elettriche, le copie fisiche vengono duplicate per gestire tempistiche differenti tra i due uffici.
    \item I contratti vengono scansionati e salvati in formato digitale. Le modifiche ed
          integrazioni vengono solo fatte sul cartaceo.
\end{itemize}

\section{Problematiche dati e software (Access vs Ad Hoc)}

\subsection{Discrepanze anagrafiche}
Esiste una disconnessione tra il database Access (gestionale
operativo/commerciale) e Ad Hoc (gestionale fiscale/amministrativo).
\begin{itemize}
    \item \textbf{Access:} Anagrafiche spesso non aggiornate , serve tenere traccia dei cambiamenti e revisioni.
    \item \textbf{Ad Hoc:} Contiene l'anagrafica fiscale corretta (Sede Legale), ma spesso le sedi operative o di spedizione sono imprecise.
    \item \textbf{Indirizzi multipli:} Necessità di distinguere sempre tre entità: Intestatario (Fiscale), Utilizzatore (Fisico), Destinatario Merce (Logistico).
\end{itemize}

\subsection{Errori commerciali e input dati}
\begin{itemize}
    \item  \textbf{Errore nella creazione:}I commerciali tendono a duplicare vecchie RA portandosi dietro errori pregressi. Attualmente, il controllo viene effettuato manualmente dal reparto spedizioni/logistica a valle del processo, causando inefficienze.
    \item \textbf{Mancazione di informazioni:} Le informazioni di trasporto quasi sempre assenti o incomplete (Incoterms, tipologia imballo, peso, volume).
\end{itemize}

\subsection{Azioni su Access eseguite da Logistica}
\begin{itemize}
    \item  \textbf{Destinazioni:}Viene corretto l'indirizzo di spedizione se errato.
    \item \textbf{Causali:} Vengono corrette le causali.
    \item \textbf{Dati di trasporto:} Vengono inseriti i dati mancanti relativi a peso, volume, imballo, Incoterms.
\end{itemize}

\section{Spedizioni Internazionali e Dogana}

\subsection{Associazione tecnico-commerciale (Packing List)}
Vi è una criticità nella creazione della \textit{Packing List} per le dogane,
specialmente extra-UE (es. Argentina).
\begin{itemize}
    \item \textbf{Macro-codici:} I commerciali vendono "macro-codici" (es. un impianto completo), ma la dogana richiede il dettaglio specifico per ogni pallet.
    \item \textbf{Problema:} La logistica deve associare manualmente le singole voci contrattuali alla distinta tecnica ai colli fisici, poiché il sistema non esplode automaticamente i macro-codici nelle loro componenti fisiche distribuite su più pallet.
    \item \textbf{Soluzione attuale:} Associare manualmente le voci dati i pallet fisici preparati dagli spedizionieri. Il problema si ripercuote alle dogane.
    \item \textbf{Problemi legati alla soluzione attuale:} L'associazione manuale fatta da UT distinta-contratto può essere errata, di coseguenza anche l'associazione fatta dagli spedizionieri pallet-distinta può essere errata.
    \item \textbf{Richiesta:} Il nuovo sistema gestionale dovrebbe permettere di collegare automaticamente le voci commerciali con quelle tecniche e facilitare il collegamneto coi pallet fisici, riducendo il rischio di errori manuali.
    \item \textbf{Gestione codici:} Necessità che Ad-hoc e il configuratore abbiano i medesimi codici articolo.
\end{itemize} 

\subsection{Gestione della vendita e modifiche:}
IL metodo attuale è focalizzato sulla vendita rapida  in fase contrattuale , ma non tiene conto delle modifiche successive che spesso avvengono in fase di progettazione e realizzazione dell'impianto. Questo comporta:
\begin{itemize}
    \item \textbf{Voci vendute:} In fase di realizzazione vengono aggiunte o modificate voci , in fase di associazione pallet-distinta tecnica-commerciale queste modifiche non sono tracciate e possono portare a discrepanze.
    \item \textbf{Modifica contratto:} Non è semplice modificare il contratto una volta firmato , questo porta a dover gestire eccezioni e variazioni in modo non strutturato.
    \item \textbf{Richiesta:} Il nuovo sistema dovrebbe guidare il commerciale in una vendita più simile alla realizzazione senza intaccare la rapidità , in questo modo si ridurrebbero le discrepanze tra voci vendute e voci realizzate.
    \item \textbf{Tracciabilità modifiche:} Il sistema dovrebbe permettere di tracciare facilmente le revisioni apportate alle voci vendute durante tutto il ciclo di vita dell'ordine.
\end{itemize}

\subsection{Vincoli normativi e alert}
Il sistema attuale manca di \textit{alert} automatici per vincoli specifici per
paese:
\begin{itemize}
    \item \textbf{Incoterms:} Necessità di definire correttamente la resa (DAP, DDP, FCA, ecc.) che determina costi e responsabilità.
    \item \textbf{ISPM 15 e Fumigazione:} Alcuni paesi (Brasile, Australia, Turchia) richiedono certificati specifici per il legno dei pallet, oltre alla marchiatura standard.
    \item \textbf{Materiali Proibiti:} Esempio di ancoranti chimici vietati in determinati paesi. Queste informazioni dovrebbero emergere automaticamente all'inserimento dell'ordine.
\end{itemize}

\section{Requisiti per il nuovo sistema gestionale}

\subsection{Storicizzazione del dato}
Il nuovo sistema deve prevedere la storicizzazione delle anagrafiche (come già
fa Ad Hoc ma non Access):
\begin{itemize}
    \item Se un cliente cambia sede o partita IVA, i documenti storici (es. DDT di anni
          passati) devono mantenere i dati validi all'epoca dell'emissione.
    \item Introduzione di \textbf{date di obsolescenza} per articoli e sedi.
\end{itemize}

\subsection{Adozione nuovo sistema:}
Per il successo del nuovo sistema, l'impostazione deve essere
\textit{top-down}:
\begin{itemize}
    \item \textbf{Regole Fisse:} Non permettere agli utenti di popolare il database liberamente. Le procedure devono essere definite a monte da un team ristretto per evitare il caos ("Più persone popolano... più casino").
    \item \textbf{Vincoli all'Ingresso:} I campi chiave (Incoterms, indirizzi, tipologia imballo) devono essere vincolati (menu a tendina) per guidare il commerciale ed evitare errori di digitazione o concettuali.
    \item \textbf{Formazione:} Necessità di formazione specifica per i commerciali sull'uso corretto del sistema e l'importanza di inserire dati accurati.
    \item \textbf{Controlli Automatici:} Implementare controlli automatici che verifichino la completezza e correttezza dei dati inseriti prima di permettere il salvataggio dell'ordine.
    \item \textbf{Guida alla Vendita:} Il sistema dovrebbe guidare il commerciale attraverso il processo di vendita, suggerendo opzioni e avvisando di potenziali errori o omissioni.
\end{itemize}

\subsection{Integrazione ideale con altri reparti}
L'obiettivo è un sistema dove:
\begin{enumerate}
    \item L'inserimento dell'ordine genera alert automatici (montaggi, permessi, vincoli
          doganali).
    \item La distinta tecnica e commerciale sono linkate (Rif CC) fin dalla
          progettazione.
    \item Il magazzino lavora su dati digitali che confluiscono direttamente nei
          documenti di trasporto senza trascrizioni manuali.
\end{enumerate}

\end{document}